\subsection*{指数把相加的参数变成复合的动作}
实指数满足参数相加、结果相乘。旋转也有同样结构：两个角度相加，对应两个旋转相继执行。因此纯虚指数与单位圆之间可能有联系，但结构相似还不是证明，需要检查它们确实满足同一数学定义或规律。

一条证明路线从幂级数出发：偶数次 $\mathrm i$ 的幂给实部交错项，奇数次给虚部交错项，分别组成余弦与正弦。另一条路线从微分方程出发，比较速度与初值。本书的动态章节展开后一种解释。

\subsection*{欧拉公式让三角关系变成指数运算}
有了 $\mathrm e^{\mathrm i\theta}=\cos\theta+\mathrm i\sin\theta$，两次旋转相乘就能写成一个指数。比较结果的实部和虚部，会重新得到和角公式。指数规则与三角恒等式不再是两套互不相干的记忆任务。

进一步写 $\mathrm e^{x+\mathrm i y}=\mathrm e^x\mathrm e^{\mathrm i y}$，实部参数控制缩放，虚部参数控制角度。这里“实部参数”指指数里的 $x$，不要误认为结果的实部就是 $\mathrm e^x$；结果实部还要乘 $\cos y$。

角度绕整圈后回到同一点，因此复指数不单射。下一步取对数时必须面对这个事实，不能简单把指数符号取消。一个好公式既能简化计算，也会暴露哪些旧习惯需要修正。
